\documentclass[preprint]{revtex4-2}
\usepackage{amsmath}
\usepackage{graphicx}
\usepackage{comment}
\graphicspath{{figures/}}
\newcommand{\aap}{Astron. Ap.}
\newcommand{\apjs}{Astrophys. J. Suppl.}
\newcommand{\gca}{Geochim. Cosmochim. Acta}
\newcommand{\epsl}{Earth Planet. Sci. Lett.}

\begin{document}

\title{NRLEE Nucleosynthesis\footnote{Supplemental material for the paper
{\it Constraints on the building blocks of the Solar System: Insights from nucleosynthetic isotope anomalies}, GCA (2021), authors: K. Bermingham,
B. S. Meyer, and K. Mezger, submitted}}
\author{Bradley S. Meyer}
\affiliation{Department of Physics and Astronomy, Clemson University, Clemson,
SC 29634, USA}
\author{Katherine R. Bermingham}
\affiliation{ Department of Earth and Planetary Sciences, Rutgers University, Piscataway, NJ 08854, USA}
\author{Klaus Mezger}
\affiliation{Institut f\"ur Geologie, Baltzerstrasse 1+3 3012 Bern, Schweiz}

\begin{abstract}
A simple, one-dimensional model of the explosion of a dense white dwarf star
is presented to illustrate the nucleosynthesis that could occur in
such events.
The initial composition of the white dwarf star
was taken to be solar abundances evolved through helium burning to model
the nuclear evolution to a carbon and oxygen white dwarf star composition.
This composition is then
subjected to an exponential distribution of neutron exposures with average
exposure $\tau_0 = 0.3\ mb^{-1}$ to model the s processing that would occur
during helium shell burning.
Single-zone nuclear reaction network calculations
were performed on that initial composition over
range of initial densities that were then taken to expand
and cool exponentially.  The calculations included electron-capture reactions
using a simple but realistic parameterization of those rates.  The results
of those calculations were converted into a white dwarf model by applying
solutions to the Lane-Emden equation to map initial densities into interior
mass coordinates in the star.  Mass fractions of resulting
cosmochemically interesting species are presented as a function of interior
mass.  The results show that in the inner, dense regions of the white dwarf star
the initial carbon and oxygen present is converted into neutron-rich iron-group
species like $^{48}$Ca and $^{50}$Ti.  In the outer, less dense regions,
the initial carbon and oxygen are not completely destroyed and a neutron burst
can produce interesting isotopes like $^{96}$Zr.  In between these layers,
oxygen burning produces substantial overabundances of $^{46}$Ti and $^{84}$Sr.
Readers may explore the results of this model in more detail via Jupyter
notebooks available at \url{http://webnucleo.org}.
\end{abstract}

\maketitle

\section{Introduction}\label{sec:intro}

Correlated isotopic anomalies in a variety of elements, but especially
in the neutron-rich isotopes of iron-group species like $^{48}$Ca, in planetary
materials appear to call for co-production of those species in a special
astrophysical environment.  This environment is not yet well-characterized,
but, 
because the nucleosynthesis of $^{48}$Ca
occurs in freezeouts from low-entropy, neutron-rich quasistatistical equilibrium
\cite{1996ApJ...462..825M}, it may generically be termed a NRLEE
(a Neutron-Rich, Low-Entropy Ejector) \cite{Bermingham2021}.  A plausible
NRLEE is a thermonuclear electron-capture supernova, which is the explosion
of a dense white dwarf star \cite{2019A&A...622A..74J}.  Since such models
are not yet well studied, it is useful to present some details of
nucleosynthesis in a NRLEE via a simple model.  That is the purpose of the
present brief paper.

\section{Initial Model} \label{sec:initial}

The initial composition of the NRLEE model was taken to be the ashes of
stellar helium burning enhanced by s-process nucleosynthesis, as would be
expected to occur in helium burning shells in the star during its
Asymptotic-Giant Branch phase.  The Webnucleo {\it mainline-nucleosynthesis}
project (\url{https://github.com/mbradle/mainline-nucleosynthesis}) was
used to compute the nucleosynthesis in the main stellar burning phases of
a star starting from initial Solar composition \cite{2003ApJ...591.1220L}.
The present model employed the default project conditions for hydrogen burning
(constant temperature $T = 2.5\times 10^7$ K, mass density $\rho = 15$ g/cc,
and a duration of 10 Myr years) and helium burning
($T = 2.5\times 10^8$K, $\rho = 1500$ g/cc, and a duration of 10,000 years).
All calculations used the Webnucleo Reaclib V2.2 XML data
(\url{https://osf.io/5cyg7/}) created from the JINA Reaclib database
(\url{https://reaclib.jinaweb.org}).
The hydrogen burning converted $^1$H into $^4$He and the initial C, N, and O
abundances predominantly into $^{14}$N by CNO cycling.  The helium burning
converted $^4$He into $^{12}$C and $^{16}$O and the $^{14}$N into $^{22}$Ne.
Some s processing occurred by release of neutrons due to reactions like
$^{22}$Ne$(\alpha,n)^{25}$Mg during helium burning.  The resulting abundances
per nucleon as a function of mass number are shown in
Fig. \ref{fig:ys_fig} in the curve labeled {\it He burning}.

During helium shell burning in the asymptotic-giant branch phase,
the material experiences s-process nucleosynthesis in a series of
pulses.  Each pulse subjects the matter to a flux of neutrons characterized
by the neutron exposure $\tau$, which is the number of neutrons passing
through an area over some time interval and is conveniently expressed in
inverse millibarns ($mb^{-1}$).  Over the course of many shell pulses, the
matter is subjected to a distribution of neutron exposures $\rho(\tau)$.
In the present model, the Webnucleo {\it simple\_s\_process} project
(\url{https://github.com/mbradle/simple_s_process}) was used to expose
the helium burning composition computed as described above to a range of
neutron exposures.  Specifically, the calculation held the initial
helium burning composition at constant temperature $T = 2.5 \times 10^8$ K,
mass density $\rho = 1000$ g/cc, and neutron number density $1.e8$ per cc
for $10^5$ years.  The calculation kept track of the neutron exposure at
each step.  The resulting abundances were then averaged over a distribution
of exposures
\begin{equation}
\rho(\tau) = \frac{1}{\tau_0} e^{-\tau/\tau_0}
\label{eq:rho_exposure}
\end{equation}
with an average exposure $\tau_0 = 0.3\ mb^{-1}$.  Such an average exposure
gives abundances that approximate well the Solar System's s-process
abundances (e.g., \cite{1974ApJ...193..397C}).  The resulting abundances
are also shown in Fig. \ref{fig:ys_fig}.  The s processing significantly
enhanced the abundance of nuclei above the iron peak.  Alpha decays of
bismuth and polonium isotopes at the end point of the s process contributed
to the fairly high abundance of $^4$He.

\section{Explosive Nucleosynthesis} \label{sec:expl}

The nucleosynthesis in the exploding white dwarf star was modeled with
the Webnucleo {\it NRLEE-Nucleosynthesis} project
(\url{https://github.com/mbradle/NRLEE-Nucleosynthesis}).  The initial
composition was that
from the distribution of neutron exposures computed as in \S \ref{sec:initial},
and a series of single-zone network calculations at initial mass densities
$\rho_0$ ranging from $10^6$ g/cc to $10^{10}$ g/cc were performed.
Each initial density
corresponds to a zone in the star.  The initial temperature
$T_9 = T / 10^9$ K
for each zone was taken to be
\begin{equation}
T_9 = 14 \left(\frac{\rho_0}{10^{10}\ \rm{g/cc}}\right)^{1/3}
\label{eq:t9_0}
\end{equation}
Fig. \ref{fig:t9_vs_rho} shows the initial temperature as a function of
initial mass density.  The matter was taken to expand exponentially on
a timescale $\tau$ such that the mass density $\rho(t)$ as a function of
time behaved as
\begin{equation}
\rho(t) = \rho_0 e^{-t / \tau}
\label{eq:rho_t}
\end{equation}
The temperature was taken to behave a $\rho \propto T^3$ and the expansions
were followed for 100 seconds.

For the purposes of understanding the evolving neutron richness of the
abundances during the explosive nucleosynthesis, the Reaclib V2.2
reaction rates were supplemented by a simple but realistic
weak-rate parameterization \cite{2010A&A...522A..25A}.  Those weak rates
included electron-capture rates on all nuclei if energetically favorable.
The electron-capture rates are fastest at high density where the density
of free electrons is high.  Fig. \ref{fig:ye_vs_time} shows the electron
fraction $Y_e$ in the zone with initial mass density $7.59 \times 10^9$ g/cc.
The electron fraction is the net number of electrons per nucleon in the system.
For charge neutral matter, this is then the fraction of all nucleons that
are protons.  In this zone, the matter starts with nearly equal numbers of
neutrons and protons.  
The density in this zone
is sufficiently high the significant electron capture occurs during the
expansion and drives down $Y_e$.  The slower the expansion, the greater
the electron capture.  For the slower expansion, $Y_e$ drops to a minimum
value and then, owing to the shifting abundances, beta-decay reactions
occur and cause an increase of $Y_e$ again.

Fig. \ref{fig:ye_vs_rho} shows the final $Y_e$ in the different zones
(indicated by their initial density).
For the lower density zones, little electron capture
occurs, and $Y_e$ is unchanged from the initial value, which is slightly
below $Y_e = 0.5$ owing to beta decays during hydrogen and helium burning
and s processing.  For the higher
density zones, significant electron capture occurs and leaves the
material significantly neutron rich.  Slower expansion allows for more
electron capture and, hence, more neutron-rich matter.
It is worth noting that the $Y_e$ of $^{48}$Ca is 20/48 = 0.42, so the
$Y_e$ values reached in the highest density zones reach values near those
to produce this isotope.

Fig. \ref{fig:c_o_ne_mg_si_vs_rho} shows the final mass fractions (grams 
of the species per total gram of the material) of
$^{12}$C, $^{16}$O, $^{20}$Ne, $^{24}$Mg, and $^{28}$Si in the different
zones.
For low initial density, the peak temperature is not
high enough to burn much of the initial carbon and oxygen except to 
convert a small amount of the oxygen to neon.  At higher
density, the carbon and even some oxygen burn during explosive carbon
burning, which produces $^{20}$Ne and $^{24}$Mg.  At higher initial density
and peak temperature, the burning proceeds through explosive carbon and
neon burning.  At still higher initial density and peak temperature,
the burning progresses through oxygen burning.  This leaves an abundance
of $^{28}$Si and heavier isotopes.  Finally, at still higher temperatures
and density, the burning proceeds all the way through silicon burning
and produces the iron-group isotopes.

In the higher initial density zones, the temperature and density are high
enough to burn matter to iron-group species.  Fig. \ref{fig:inner_x_vs_time}
shows the mass fractions of several species as a function of time
in the zone with initial mass density $7.59 \times 10^9$ g/cc for the
$\tau = 0.1$ s expansion.
Due to the high temperature and density, the
initial carbon and oxygen in the zone transform within about $10^{-12}$
seconds into a nuclear statistical equilibrium dominated by neutrons,
protons, and $^4$He nuclei.  Nevertheless, heavier species are also
present in relatively large abundance due to the low entropy.
Around $10^{-3}$ seconds into the explosion, the electron-capture
reactions start decreasing the $Y_e$, which favors the build up of
neutron-rich species, like $^{50}$Ti.  The density (and temperature)
start to decline rapidly starting around $10^{-2}$ seconds; hence,
it is favorable for the light particles to assemble into heavier species.
At late time, the abundances freeze out except for the decay of
radioactive species.

In the lower-density zones, leftover $^{22}$Ne from helium burning can provide
neutrons during the explosive nucleosynthesis.  Fig. \ref{fig:n_n_vs_time}
shows the neutron number density as a function of time in the zone with
initial mass density $7.59 \times 10^6$ g/cc for the $\tau = 0.1$ s expansion.
A number of reactions, but particularly
$^{22}$Ne$(\alpha, n)^{25}$Mg release neutrons during the explosion.
This leads to a sudden, large flux of neutrons.

The burst of neutrons in lower-density zones can significantly modify the
initial abundances during the explosion.  Fig. \ref{fig:zr_vs_time}
shows the evolution of the zirconium mass fractions
in the zone with initial mass density $7.59 \times 10^6$ g/cc for the
$\tau = 0.1$ s expansion.
The sudden burst of neutrons causes the zirconium
isotopes to capture neutrons rapidly.  The initially relatively small
abundance of $^{96}$Zr rapidly builds up from $^{90}$Zr and the other
zirconium isotopes \cite{2011ApJ...739...93T}.
It is worth noting that the s-process nucleosynthesis
that occurred prior to the explosion significantly enhanced the zirconium
abundances (see Fig. \ref{fig:ys_fig}).

The final mass fractions of species as a function of initial density vary
owing to the different conditions experienced during the explosive
nucleosynthesis.  Fig. \ref{fig:ti_vs_rho} shows the final titanium
mass fractions as a function of the initial mass density for the $\tau = 0.1$
s expansions.
The high mass fraction
of $^{46}$Ti is noteworthy near initial density of $10^8$ g/cc.  It arises
because the explosive oxygen burning in that region occurs in slightly
neutron-rich conditions, which favors $^{46}$Ti over $^{44}$Ti in the
nuclear flows \cite{1973ApJS...26..231W}.
The figure includes the abundant $^{48}$Cr present
because that species decays to $^{48}$Ti.

\section{White Dwarf Model} \label{sec:polytrope}

The calculations in \S \ref{sec:expl} are useful in delineating the 
varying yields that could result in conditions relevant to NRLEE
nucleosynthesis.  It is helpful, however, for understanding the physical
setting of this nucleosynthesis to express those results in the context
of a white dwarf star model.  It is possible to do this
through the use of a polytrope model
(e. g., \cite{1984psen.book.....C}).
For a spherically symmetric star, the mass-continuity equation reads
\begin{equation}
\frac{dM_r}{dr} = 4 \pi r^2 \rho
\label{eq:continuity}
\end{equation}
where $r$ is the radial coordinate, $\rho$ is the mass density, and
$M_r$ is the enclosed mass (the mass contained within a sphere of radius $r$).
For a spherically symmetric star in hydrostatic equilibrium, the
pressure gradient force balances gravity such that
\begin{equation}
\frac{dP}{dr} = -\frac{G M_r\rho}{r^2}
\label{eq:hydro}
\end{equation}
where $P$ is the pressure and $G$ is Newton's constant.

To construct a polytrope, one uses a polytropic equation of state,
which is given by
\begin{equation}
P = K \rho^{(n+1)/n}
\label{eq:polytrope}
\end{equation}
with $K$ a constant.  The density can be defined in terms of a function
$\theta(\xi)$ such that
\begin{equation}
\rho = \rho_c \theta^n
\label{eq:rho}
\end{equation}
where $\rho_c$ is the central density and $n$ is the polytropic index.
Eqs. (\ref{eq:continuity}), (\ref{eq:hydro}), (\ref{eq:polytrope}),
and (\ref{eq:rho}) can be combined to yield the Lane-Emden equation
\begin{equation}
\frac{1}{\xi^2} \frac{d}{d\xi}\left(\xi^2\frac{d\theta}{d\xi}\right)
+ \theta^n = 0
\label{eq:lane-emden}
\end{equation}
where
\begin{equation}
r = a \xi
\label{eq:r_xi}
\end{equation}
and
\begin{equation}
a^2 = \frac{(n + 1)K\rho_c^{\frac{1}{n} - 1}}{4\pi G}
\label{eq:a}
\end{equation}
With a choice of $n$ (allowed values are $0 \leq n \leq 5$), one solves
Eq. (\ref{eq:lane-emden}) for $\theta(\xi)$.

The total radius of the polytrope model is given by $R = a \xi_1$,
where $\xi_1$ is the first root of $\theta$ such that $\theta(\xi_1) = 0$.
The mass of the polytrope model out to coordinate $\xi$ is found by
using Eq. (\ref{eq:lane-emden}) to
integrate Eq. (\ref{eq:continuity}) to obtain
\begin{equation}
M(\xi) = -4 \pi a^3\rho_c \xi^2 \frac{d\theta}{d\xi}
\label{eq:m_xi}
\end{equation}
with the total mass given by ${\cal M} = M(\xi_1)$.  Choice of $n$,
a total mass $\cal{M}$, and central density $\rho_c$ then determines $a$ and
$r$ as a function of $\xi$, and, thus, the density
and mass as a function of $r$.  In this way, one can obtain the enclosed
mass coordinate $M_r = M(r)$ as a function of density $\rho$.
Since a NRLEE is expected to be the explosion of a dense white dwarf star
with mass near the the Chandrasekhar mass, appropriate choices are
${\cal M} = 1.4\ M_\odot$ ($M_\odot$ is one Solar mass) and $n = 3$ (since
the Chandrasekhar limit is when nearly all the degenerate electrons 
providing the pressure support are relativistic).  If one then chooses
a central density $\rho_c = 10^{10}$ g/cc, a reasonable value for such a
dense white dwarf star, the run of density with enclosed mass is that shown
in Fig. \ref{fig:rho_vs_mr}.

With the mapping between the density and enclosed mass now available, one
can simply convert other quantities from dependence on initial density to
enclosed mass.  Fig. \ref{fig:t9_vs_mr} shows the peak temperature in the
explosion as a function of enclosed mass.  The highest peak temperature is
reached at the center of the exploding star and decreases outward.  It is
important to note that, in a real star, the explosion is a dynamical
event such that release of energy
in one part of the star causes the
various stellar layers to respond by heating up and expanding.  This means
that various stellar zones do not reach their peak temperatures at the same
time.  The present model does not capture those complex dynamics
since the layers are taken to heat up and expand independently.  Nevertheless,
since the explosion occurs quickly compared to the timescale to significantly
change the structure of the star, the present treatment gives a fair
approximation of the real explosion, at least in one dimension.

Fig. \ref{fig:ye_vs_mr} shows the final $Y_e$ immediately after the
explosion as a function of the enclosed mass coordinate.  This is a mapping
of Fig. \ref{fig:ye_vs_rho} from initial density to enclosed mass.
As expected, the inner regions of the white dwarf are densest and experience
the most neutron capture.  Also, the slower the explosion (the slower the
expansion), the more neutron-rich the matter becomes.

Fig. \ref{fig:ti_c_o_vs_mr} shows the mass fractions of the titanium 
and $^{12}$C and $^{16}$O immediately after the explosion.  The figure is
for the $\tau = 0.1$ s expansion, as are all subsequent figures plotting
a quantity vs. enclosed mass.  The unburned
or slightly burned carbon and oxygen are in the zones that attained the
lowest peak temperatures.  In the present model, these are in the outer
few tenths of a Solar mass of the exploded star.
The inner layers of the star are dominated by
neutron-rich $^{50}$Ti.  The $^{46}$Ti-rich layers are in the oxygen-rich,
outer parts of the star.

For understanding the relevance of the nucleosynthesis in the NRLEE for
cosmochemistry, it is helpful to plot the overabundances instead of mass
fraction.  The overabundance of a species is its mass fraction in a
nucleosynthetic environment relative to its mass fraction in the Solar System.
Thus, since the mass fraction can be labeled $X$, the overabundance is
labeled $X / X_\odot$.  Fig. \ref{fig:ti_over_vs_mr} shows the titanium
overabundances in the simple polytrope model.
Zones with huge overabundances of $^{50}$Ti are present in the
inner parts of the exploded star.   Zones with significant overabundances
of $^{46}$Ti are present in the outer, oxygen-rich regions of the star.
The $^{48}$Cr overabundance shown
is the overabundance in $^{48}$Ti that would be present after all the
$^{48}$Cr decayed via the chain $^{48}\rm{Cr} \to ^{48}\rm{V} \to ^{48}\rm{Ti}$.
The half life of $^{48}$Cr is 21.6 hours and that of $^{48}$V is 16 days.
The ejecta may condense into dust prior to the time for this decay chain
to run to completion, so the overabundance shown is the maximum one would
expect in a condensate for the present model.
It is clear that $^{46}$Ti and $^{50}$Ti overabundances dominate in this
simple model, which, if they co-condense in the NRLEE outflow, would result
in carriers with correlated enrichments in these isotopes that could then
be the source of correlated anomalies in planetary materials
\cite{2009Sci...324..374T}.

Fig. \ref{fig:ca_over_vs_mr} shows the overabundances of the calcium isotopes
in the simple polytrope model.
Zones with huge overabundances of $^{48}$Ca are present in the
inner parts of the exploded star.   Zones with significant overabundances
of $^{42}$Ca and $^{46}$Ca  are present in the outer, oxygen-rich regions of
the star.  If neutron-rich material from the inner part of the exploded
star indeed mixed out to the oxygen-rich layers to form oxides, one could
naturally expect correlated $^{48}$Ca and $^{50}$Ti, especially if the
isotopes condensed in perovskite.

Fig. \ref{fig:cr_over_vs_mr} shows the overabundances of the chromium isotopes
in the simple polytrope model.
Zones with huge overabundances of $^{54}$Ca are present in the
inner parts of the exploded star.  Chromium-50 dominates the overabundance
over a large mass range outside the $^{54}$Cr-rich region.  Of course,
one would expect the highly enriched $^{54}$Cr to accompany the 
abundant $^{48}$Ca and $^{50}$Ti as the inner layers mixed to the outer,
oxygen-rich layers.

Fig. \ref{fig:fe_over_vs_mr} shows the iron isotopes in the model.
Inside of the oxygen-rich outer layers, $^{54}$Fe is
highly overabundant.  Inside of these layers, $^{58}$Fe is highly
overabundant.

Fig. \ref{fig:ni_over_vs_mr} shows the nickel overabundances in the
simple polytrope model.
Isotopes of increasing neutron richness come to dominate
the overabundances with increasing depth in the exploded star.  Closest
to the oxygen-rich outer layers is the $^{58}$Ni-rich region.  The proximity
of this region to the oxygen might lead to nickel bearing condensates
preferentially enriched in $^{58}$Ni over the other nickel isotopes.
The figure also shows the $^{60}$Ni overabundance that would result from
$^{60}$Fe in the innermost layers after that isotope decayed.

The co-existing dominant overabundances of $^{54}$Fe and $^{58}$Ni
in the interior mass coordinate range $\sim 0.7 - 1.2\ M_\odot$ are
striking in light of the correlated excesses in these isotopes
in CC meteorites
\cite{Hopp2021}.  These isotopes are produced at densities that reach
$Y_e \approx 0.48$ during the expansion.  This is because
$Y_e(^{54}{\rm Fe}) = 26 / 54 = 0.4815$ 
and
$Y_e(^{58}{\rm Ni}) = 28 / 58 = 0.4828$; thus, at the high temperature
and density achieved in the explosion, as the fast nuclear reactions
drive the nuclear abundance distribution into nuclear statistical equilibrium
(NSE)
and the weak reactions drive the $Y_e$ to $\sim 0.48$. In the equilibrium,
the nucleons arrange
themselves into the nuclei with strongest binding energy per nucleon for
the given $Y_e$.  The Fe and Ni isotopes dominate the binding energy per nucleon
and $^{54}$Fe and $^{58}$Ni have the right electron fraction.  The binding
energy per nucleon of $^{54}$Fe is 8.7362 MeV per nucleon while that of
$^{58}$Ni is 8.7319 MeV per nucleon.  Fig. \ref{fig:fe_ni_vs_mr}
shows the final mass fractions of the iron and nickel isotopes in the
polytrope model.  The slight nuclear binding energy
advantage of $^{54}$Fe allows it to be more abundant than $^{58}$Ni in
the freezeout from the NSE.  The system includes a small amount of other
Fe and Ni isotopes with other $Y_e$ values
to accommodate the total electron fraction.

Fig. \ref{fig:sr_over_vs_mr_linear} shows the strontium overabundances
in the simple polytrope model.
Only the outermost layers of the exploded star are shown
since the burning destroys all strontium inside of these layers.  The high
$^{84}$Sr overabundance over a narrow region arises from gamma processing.

Fig. \ref{fig:zr_over_vs_mr_linear} shows the zirconium overabundances
in the simple polytrope model.
Only the outermost layers of the exploded star are shown
since the burning destroys all zirconium inside of these layers.  The
high $^{96}$Zr in the outermost layers
arises from the neutron burst processing there
(see Fig. \ref{fig:zr_vs_time}).

Fig. \ref{fig:mo_over_vs_mr_linear} shows the molybdenum overabundances
in the simple polytrope model.
Only the outermost layers of the exploded star are shown
since the burning destroys all molybdenum inside of these layers.
The innermost layers are enriched in the p-process isotopes $^{92,94}$Mo
due to the gamma processing occurring there.  Outside of these layers,
other Mo isotopes dominate.  Interestingly, nowhere is $^{95}$Mo the most
overabundant species.  On the other hand, the figure shows that
$^{95}$Zr ($\tau_{1/2} = 64.0$ days) and $^{95}$Nb ($\tau_{1/2} = 35.0$ days)
can give a significant overabundance
in $^{95}$Mo in NRLEE dust after {\it in situ} decay.  It should be
noted that the approximate weak rates used for the present models
\cite{2010A&A...522A..25A}
overestimate the decay rate of $^{95}$Zr, so the overabundance from
$^{95}$Zr is probably lower than is realistic while that from
$^{95}$Nb is probably higher than is realistic.  Summing the two overabundances
in the figure probably gives a good estimate of the resulting $^{95}$Mo
overabundance if Zr and Nb condense in NRLEE dust.
Though not shown, $^{96}$Nb would contribute insignificantly to the 
overabundance of $^{96}$Mo.
In any case, it is evident that NRLEE dust could carry enrichments
in $^{94}$Mo and $^{95}$Mo relative to $^{96}$Mo, which could then be
the source of correlated anomalies in $^{94}$Mo and $^{95}$Mo
seen in planetary materials \cite{2019NatAs...3..736B}.
This possibility merits further quantitative work.

\section{Jupyter Notebooks} \label{sec:jupyter}

Readers who wish to study the above results in 
more detail are encouraged to explore the Jupyter notebooks accompanying
the Webnucleo projects used to construct the simple NRLEE model presented
here.  The notebooks allow a user to download and graph data interactively.
Input parameters can be changed, and figures and tables saved.  The
notebooks run on Google Colab (\url{https://colab.research.google.com)} or
Binder (\url{https://mybinder.org}).  They can also be downloaded and run
on a local computer.

The notebooks relevant for the model presented here are:
\begin{itemize}
\item Mainline-nucleosynthesis: \url{https://github.com/mbradle/mainline-nucleosynthesis}
\item simple\_s\_process: \url{https://github.com/mbradle/simple_s_process}
\item NRLEE-Nucleosynthesis: \url{https://github.com/mbradle/NRLEE-Nucleosynthesis}
\end{itemize}
All of these notebooks are part of the Webnucleo Project
(\url{http://webnucleo.org}).

\section{Conclusion}

A simple, one-dimensional model of a NRLEE (Neutron-Rich Low-Entropy
Ejector) has been presented.  The model was constructed from independent
single-zone network calculations that were then stitched together into
an exploding white-dwarf star model through the use of an $n = 3$, 
${\cal M} = 1.4\ M_\odot$ polytrope configuration.
Fig. \ref{fig:nrlee-anatomy} summarizes the NRLEE nucleosynthesis.
The innermost, densest zones reach the highest temperature and density
during the explosion.  They experience significantly fast electron-capture
reactions during the explosion such that the matter becomes neutron rich
to produce $^{48}$Ca, $^{50}$Ti, and $^{54}$Cr.  Outside those regions,
less neutron-rich, intermediate mass isotopes are produced.  In even
lower-density zones, initial oxygen does not completely burn.  In these
zones, abundant $^{46}$Ti and $^{84}$Sr are produced.  Finally, in the outer
zones, unburned carbon remains along with isotopes like $^{96}$Zr produced
by neutron-burst nucleosynthesis.

If the layers of this simple model mix and then condense dust grains, those
grains may be important contributors to the isotopic anomalies in the
early Solar System \cite{Bermingham2021}.  For example, the model ejecta
are significantly enriched in $^{46}$Ti and $^{50}$Ti.  If the layers in the
ejecta carrying these enrichments mix and then condense, the resulting NRLEE
dust could be the source of correlated anomalies in these two isotopes in
Solar System planetary samples.

The present model is obviously extremely simple.  The initial composition
for the explosion calculations is computed from constant temperature and
density single-zone calculations and then modified by a schematic treatment
of the s process that would likely occur during the pre-explosion star's
shell helium burning phase.  The explosion calculations are independent
single-zone calculations with initial temperature derived from a simple
density-temperature scaling and the subsequent dynamics given by a simple
exponential expansion.  The results are then mapped to a one-dimensional
white-dwarf star
model through use of a simple polytrope configuration.  Nevertheless, the
model gives something of the flavor of what more realistic,
multi-dimensional  models are likely
to produce in the future.  Such detailed models are eagerly anticipated!

\bibliography{nrlee-nucleosynthesis}

\begin{figure}
\includegraphics{ys_fig}
\caption{The abundances as a function of mass number
after helium burning in the model (He burning) and after those abundances are
subjected to an exponential distribution of neutron exposures with average
exposure of 0.3 $mb^{-1}$.}
\label{fig:ys_fig}
\end{figure}

\begin{figure}
\includegraphics{t9_vs_rho}
\caption{The peak temperature in the model calculations as a function of
the initial density. The relationship between the initial density $\rho$
and peak temperature was taken to be $\rho \propto T_9^3$.}
\label{fig:t9_vs_rho}
\end{figure}

\begin{figure}
\includegraphics{ye_vs_time}
\caption{The evolution of the electron fraction $Y_e$ as a function of time
in the zone with initial mass density $7.59 \times 10^9$ g/cc for the
indicated density expansion timescales.
}
\label{fig:ye_vs_time}
\end{figure}

\begin{figure}
\includegraphics{ye_vs_rho}
\caption{The final electron fraction $Y_e$ (100 seconds after the
explosion) as a function of 
initial density for the indicated density expansion timescales.
}
\label{fig:ye_vs_rho}
\end{figure}

\begin{figure}
\includegraphics{c_o_ne_mg_si_vs_rho}
\caption{The final $^{12}$C, $^{16}$O, $^{20}$Ne, $^{24}$Mg, and $^{28}$Si
mass fractions (100 s after the
explosion) as a function of 
initial density.
}
\label{fig:c_o_ne_mg_si_vs_rho}
\end{figure}

\begin{figure}
\includegraphics{inner_x_vs_time}
\caption{
The evolution of the mass fractions of several species as a function of time
in the zone with initial mass density $7.59 \times 10^9$ g/cc for the
$\tau = 0.1$ s expansion.
}
\label{fig:inner_x_vs_time}
\end{figure}

\begin{figure}
\includegraphics{n_n_vs_time}
\caption{
The evolution of the neutron number density (in grams per cubic centimeter)
in the zone with initial mass density $7.59 \times 10^6$ g/cc for the
$\tau = 0.1$ s expansion.
}
\label{fig:n_n_vs_time}
\end{figure}

\begin{figure}
\includegraphics{zr_vs_time}
\caption{
The evolution of the zirconium mass fractions
in the zone with initial mass density $7.59 \times 10^6$ g/cc for the
$\tau = 0.1$ s expansion.
}
\label{fig:zr_vs_time}
\end{figure}

\begin{figure}
\includegraphics{ti_vs_rho}
\caption{
The final titanium mass fractions (100 s after the explosion) as a function
of initial density for the $\tau = 0.1$ s expansion.
}
\label{fig:ti_vs_rho}
\end{figure}

\begin{figure}
\includegraphics{rho_vs_mr}
\caption{
The density as a function of interior mass coordinate $M_r$ for a simple
$n = 3$, $M = 1.4\ M_\odot$ polytrope.
}
\label{fig:rho_vs_mr}
\end{figure}

\begin{figure}
\includegraphics{t9_vs_mr}
\caption{
The peak temperature
as a function of interior mass coordinate $M_r$ for a simple
$n = 3$, $M = 1.4\ M_\odot$ polytrope.
}
\label{fig:t9_vs_mr}
\end{figure}

\begin{figure}
\includegraphics{ye_vs_mr}
\caption{The final electron fraction $Y_e$ (100 seconds after the
explosion) as a function of 
enclosed mass for the indicated density expansion timescales.
}
\label{fig:ye_vs_mr}
\end{figure}

\begin{figure}
\includegraphics{ti_c_o_vs_mr}
\caption{
The titanium and $^{12}$C and $^{16}$O mass fractions in the simple
polytrope model.
}
\label{fig:ti_c_o_vs_mr}
\end{figure}

\begin{figure}
\includegraphics{ti_over_vs_mr}
\caption{
The titanium overabundances in the simple
polytrope model.
}
\label{fig:ti_over_vs_mr}
\end{figure}

\begin{figure}
\includegraphics{ca_over_vs_mr}
\caption{
The calcium overabundances in the simple
polytrope model.
}
\label{fig:ca_over_vs_mr}
\end{figure}

\begin{figure}
\includegraphics{cr_over_vs_mr}
\caption{
The chromium overabundances in the simple
polytrope model.
}
\label{fig:cr_over_vs_mr}
\end{figure}

\begin{figure}
\includegraphics{fe_over_vs_mr}
\caption{
The iron overabundances in the simple
polytrope model.
}
\label{fig:fe_over_vs_mr}
\end{figure}

\begin{figure}
\includegraphics{ni_over_vs_mr}
\caption{
The nickel overabundances in the simple
polytrope model.
}
\label{fig:ni_over_vs_mr}
\end{figure}

\begin{figure}
\includegraphics{fe_ni_vs_mr}
\caption{
The final iron and nickel  in the simple
polytrope model (100 s after the explosion).
}
\label{fig:fe_ni_vs_mr}
\end{figure}

\begin{figure}
\includegraphics{sr_over_vs_mr_linear}
\caption{
The strontium overabundances in the simple
polytrope model.
}
\label{fig:sr_over_vs_mr_linear}
\end{figure}

\begin{figure}
\includegraphics{zr_over_vs_mr_linear}
\caption{
The zirconium overabundances in the simple
polytrope model.
}
\label{fig:zr_over_vs_mr_linear}
\end{figure}

\begin{figure}
\includegraphics{mo_over_vs_mr_linear}
\caption{
The molybdenum overabundances in the simple polytrope model.  The curves
labeled $^{95}$Zr and $^{95}$Nb are the overabundances of $^{95}$Mo that
would result from decay of parent $^{95}$Zr and $^{95}$Nb, respectively.
}
\label{fig:mo_over_vs_mr_linear}
\end{figure}

\begin{figure}
\includegraphics[width=\textwidth]{nrlee-anatomy}
\caption{
Possible anatomy of a NRLEE.
}
\label{fig:nrlee-anatomy}
\end{figure}

\end{document}
